% F085｜八类各有例题，优先用特殊值求解，避免先猜解析式。
\FReserve{4cm}
\begin{examplebox}
\textbf{\textcolor{CExample}{例 1（指数型）｜已知 $f(1)$，不求解析式}}

在全体实数上有 $f(x+y)=f(x)f(y)$，且 $f(1)=3$。求 $f(4)$。
\[
f(4)=f(1+1+1+1)=f(1)^4=\boxed{81}.
\]
\textbf{最快：}看到加法就把 $4$ 拆成四个 $1$；不必猜 $f(x)=3^x$。回验 $f(2)=9$，再有 $f(4)=f(2)^2=81$。
% F085｜真实数学图；代表函数不构成抽象解的唯一性证明。
\begin{center}
\begin{tikzpicture}
\begin{axis}[width=12.4cm,height=5.1cm,axis lines=middle,
 xmin=-0.8,xmax=4.5,ymin=-4,ymax=100,
 xtick={0,1,2,3,4},xticklabels=\empty,ytick={0,9,27,81},xlabel={$t$},ylabel={$f(t)$},
 title={例 1 的代表曲线 $3^t$（不声称唯一）},title style={font=\small,yshift=3pt},
 tick label style={font=\small},label style={font=\small},
 ylabel style={at={(axis description cs:0,1)},anchor=south east,xshift=-18pt},
 clip=false,grid=major,grid style={black!8},enlargelimits=false,
 every axis plot/.append style={thick},samples=161]
\addplot[CExplanation,very thick,domain=-.5:4.1] {3^x};
\draw[dashed,black!35] (axis cs:1,0) |- (axis cs:0,3);
\draw[dashed,black!35] (axis cs:2,0) |- (axis cs:0,9);
\draw[dashed,black!35] (axis cs:3,0) |- (axis cs:0,27);
\draw[dashed,black!35] (axis cs:4,0) |- (axis cs:0,81);
\addplot[only marks,mark=*,mark size=2pt,orange!80!black] coordinates {(1,3) (2,9) (3,27) (4,81)};

\node[font=\small,anchor=north,yshift=-3pt] at (axis cs:1,-4) {$1$};
\node[font=\small,anchor=north,yshift=-3pt] at (axis cs:2,-4) {$2$};
\node[font=\small,anchor=north,yshift=-3pt] at (axis cs:3,-4) {$3$};
\node[font=\small,anchor=north,yshift=-3pt] at (axis cs:4,-4) {$4$};

\end{axis}
\end{tikzpicture}
\par\small 输入每加 $1$，高度就乘 $f(1)=3$；四次相乘得到 $f(4)=81$。
\end{center}

\end{examplebox}

\FReserve{4cm}
\begin{examplebox}
\textbf{\textcolor{CExample}{例 2（对数型）｜乘积拆开做加法}}

在正数范围满足 $f(xy)=f(x)+f(y)$，且 $f(2)=3$。求 $f(8)$ 与 $f(1/2)$。
\[
f(8)=f(2\cdot2\cdot2)=3f(2)=\boxed{9},
\]
\[
f(1/2)=f(1)-f(2)=\boxed{-3}.
\]
\textbf{最快：}先令 $y=1$ 得 $f(1)=0$，再拆正数的积或倒数。回验 $f(2) + f(1/2)=0=f(1)$。
\end{examplebox}

\FReserve{4cm}
\begin{examplebox}
\textbf{\textcolor{CExample}{例 3（幂函数型）｜同一个 $f(2)$，结果却不同}}

在正数范围满足 $f(xy)=f(x)f(y)$，且 $f(2)=3$。求 $f(8)$ 与 $f(1/2)$。
\[
f(8)=f(2)^3=\boxed{27},\qquad
f(1/2)=\frac1{f(2)}=\boxed{\frac13}.
\]
\textbf{对照上一题：}对数模型的 $f(8)=9$，幂模型的 $f(8)=27$；区别就在等式右边是\textbf{加}还是\textbf{乘}。回验 $f(2)f(1/2)=3\cdot(1/3)=1=f(1)$。
\end{examplebox}

\FReserve{4cm}
\begin{examplebox}
\textbf{\textcolor{CExample}{例 4（正比例型）｜负数与倍数无需猜函数}}

若 $f(x+y)=f(x)+f(y)$，且 $f(2)=5$，求 $f(-6)$。
\[
f(-6)=-f(6)=-3f(2)=\boxed{-15}.
\]
\textbf{为什么能变负号？}由 $f(0)=0$ 和 $0=f(x)+f(-x)$，得 $f(-x)=-f(x)$。回验 $f(6)+f(-6)=15-15=0$。
\end{examplebox}

\FReserve{4cm}
\begin{examplebox}
\textbf{\textcolor{CExample}{例 5（余弦型）｜令 $y=x$ 就够了}}

若不恒为零的函数在实数范围内满足
\[f(x+y)+f(x-y)=2f(x)f(y),\qquad f(1)=\tfrac12,\]
求 $f(2)$。

先令 $y=0$ 得 $f(0)=1$；再令 $x=y=1$：
\[
f(2)+1=2f(1)^2=\tfrac12,\quad\boxed{f(2)=-\tfrac12}.
\]
\textbf{最快：}二倍角，只需一个数 $f(1)$ 与 $f(0)=1$。回验 $-1/2+1=2(1/2)^2$。
% F085｜真实数学图；代表函数不构成抽象解的唯一性证明。
\begin{center}
\begin{tikzpicture}
\begin{axis}[width=12.4cm,height=5.1cm,axis lines=middle,
 xmin=-0.2,xmax=6.5,ymin=-1.5,ymax=1.5,
 xtick={0,1,2,3,4,5,6},xticklabels=\empty,ytick={-1,-0.5,0,0.5,1},xlabel={$t$},ylabel={$f(t)$},
 title={例 5 的代表曲线 $\cos(\pi t/3)$},title style={font=\small,yshift=3pt},
 tick label style={font=\small},label style={font=\small},
 ylabel style={at={(axis description cs:0,1)},anchor=south east,xshift=-18pt},
 clip=false,grid=major,grid style={black!8},enlargelimits=false,
 every axis plot/.append style={thick},samples=161]
\addplot[CExplanation,very thick,domain=0:6.1] {cos(60*x)};
\draw[dashed,black!35] (axis cs:1,0) |- (axis cs:0,0.5);
\draw[dashed,black!35] (axis cs:2,0) |- (axis cs:0,-0.5);
\addplot[only marks,mark=*,mark size=2pt,orange!80!black] coordinates {(1,0.5) (2,-0.5)};

\node[font=\small,anchor=north,yshift=-3pt] at (axis cs:1,-1.5) {$1$};
\node[font=\small,anchor=north,yshift=-3pt] at (axis cs:2,-1.5) {$2$};
\node[font=\small,anchor=north,yshift=-3pt] at (axis cs:3,-1.5) {$3$};
\node[font=\small,anchor=north,yshift=-3pt] at (axis cs:4,-1.5) {$4$};
\node[font=\small,anchor=north,yshift=-3pt] at (axis cs:5,-1.5) {$5$};
\node[font=\small,anchor=north,yshift=-3pt] at (axis cs:6,-1.5) {$6$};

\end{axis}
\end{tikzpicture}
\par\small 二倍输入的高度由 $f(2)=2f(1)^2-1$ 确定；图中横轴是 $t$，不是弧度角 $\pi t/3$。
\end{center}

\end{examplebox}

\FReserve{4cm}
\begin{examplebox}
\textbf{\textcolor{CExample}{例 6（正切型）｜分母别忘了}}

函数在所用输入处满足
\[f(x+y)=\frac{f(x)+f(y)}{1-f(x)f(y)},\qquad f(1)=\tfrac12.\]
若 $1,2$ 处有定义，求 $f(2)$。
\[
f(2)=\frac{\tfrac12+\tfrac12}{1-\tfrac14}=\boxed{\tfrac43}.
\]
\textbf{检查：}分母 $1-\tfrac14=\tfrac34\ne0$；回验 $(4/3)(3/4)=1=f(1)+f(1)$。
\end{examplebox}

\FReserve{4cm}
\begin{examplebox}
\textbf{\textcolor{CExample}{例 7（正弦型）｜用中点和半差凑参数}}

已知
\[f(x+y)f(x-y)=f^2(x)-f^2(y),\quad f(3)=5,\ f(1)=2,\]
求 $f(4)f(2)$。

因为 $4=3+1$、$2=3-1$，直接代入：
\[
f(4)f(2)=f(3)^2-f(1)^2=25-4=\boxed{21}.
\]
\textbf{反过来想：}见 $f(m)f(n)$，取 $x=(m+n)/2$，$y=(m-n)/2$。回验 $21+2^2=5^2$。
% F085｜真实数学图；代表函数不构成抽象解的唯一性证明。
\begin{center}
\begin{tikzpicture}[x=1.8cm,y=1cm,>=stealth]
\draw[->] (0,0)--(5.6,0) node[right,font=\small] {输入};
\foreach \t in {1,2,3,4,5} {\draw (\t,.06)--(\t,-.06) node[below,font=\small] {$\t$};}
\fill[CExplanation] (2,0) circle(.045) (4,0) circle(.045);
\fill[orange!80!black] (3,0) circle(.045);
\draw[<->,thick,orange!80!black] (2,.55)--node[above,font=\small] {$y=1$}(3,.55);
\draw[<->,thick,orange!80!black] (3,.55)--node[above,font=\small] {$y=1$}(4,.55);
\node[font=\small,above] at (2,1.1) {$n=x-y=2$};
\node[font=\small,above] at (4,1.1) {$m=x+y=4$};
\node[font=\small,below] at (3,-.65) {中点 $x=3$，半差 $y=1$};
\end{tikzpicture}
\par\small 数轴上端点 $2,4$ 关于 $3$ 对称，所以直接代 $x=3,y=1$；无需猜测抽象函数的曲线。
\end{center}

\end{examplebox}

\FReserve{4cm}
\begin{examplebox}
\textbf{\textcolor{CExample}{例 8（余切型）｜先求 $f(2)$ 再求 $f(3)$}}

若在有关输入处满足
\[f(x+y)=\frac{f(x)f(y)-1}{f(x)+f(y)},\qquad f(1)=2,\]
求 $f(2)$、$f(3)$。
\[
f(2)=\frac{2\cdot2-1}{2+2}=\boxed{\tfrac34},
\]
\[
f(3)=\frac{\tfrac34\cdot2-1}{\tfrac34+2}=\boxed{\tfrac2{11}}.
\]
每次都应检查分母：$4\ne0$、$\frac{11}{4}\ne0$。回验 $(3/4)\cdot4=3$，$(2/11)\cdot(11/4)=1/2$，恰为两次计算的分子。
\end{examplebox}

\FReserve{4cm}
\begin{examplebox}
\textbf{\textcolor{CExample}{例 9（辨识）｜看右边是加还是乘}}

已知正数范围内 $f(2)=3$。
\begin{enumerate}[leftmargin=*,itemsep=.5em]
\item 若 $f(xy)=f(x)+f(y)$，则 $f(4)=f(2)+f(2)=\boxed{6}$。
\item 若 $f(xy)=f(x)f(y)$，则 $f(4)=f(2)^2=\boxed{9}$。
\end{enumerate}
\textbf{陷阱：}两题都是 $xy$，\textbf{真正区分它们的是等式右边}。
% F085｜真实数学图；代表函数不构成抽象解的唯一性证明。
\begin{center}
\begin{minipage}{.49\linewidth}\centering
\begin{tikzpicture}
\begin{axis}[width=6.1cm,height=4.8cm,axis lines=middle,
 xmin=-0.2,xmax=8.9,ymin=-4,ymax=30,
 xtick={0,2,4,8},xticklabels=\empty,ytick={0,3,6,9,27},xlabel={$t$},ylabel={$f(t)$},
 title={同尺度对照：$3\log_2t$},title style={font=\small,yshift=3pt},
 tick label style={font=\small},label style={font=\small},
 ylabel style={at={(axis description cs:0,1)},anchor=south east,xshift=-18pt},
 clip=false,grid=major,grid style={black!8},enlargelimits=false,
 every axis plot/.append style={thick},samples=161]
\addplot[CExplanation,very thick,domain=.5:8.1] {3*ln(x)/ln(2)};
\draw[dashed,black!35] (axis cs:2,0) |- (axis cs:0,3);
\draw[dashed,black!35] (axis cs:4,0) |- (axis cs:0,6);
\draw[dashed,black!35] (axis cs:8,0) |- (axis cs:0,9);
\addplot[only marks,mark=*,mark size=2pt,orange!80!black] coordinates {(2,3) (4,6) (8,9)};

\node[font=\small,anchor=north,yshift=-3pt] at (axis cs:2,-4) {$2$};
\node[font=\small,anchor=north,yshift=-3pt] at (axis cs:4,-4) {$4$};
\node[font=\small,anchor=north,yshift=-3pt] at (axis cs:8,-4) {$8$};

\end{axis}
\end{tikzpicture}
\end{minipage}\hfill
\begin{minipage}{.49\linewidth}\centering
\begin{tikzpicture}
\begin{axis}[width=6.1cm,height=4.8cm,axis lines=middle,
 xmin=-0.2,xmax=8.9,ymin=-4,ymax=30,
 xtick={0,2,4,8},xticklabels=\empty,ytick={0,3,6,9,27},xlabel={$t$},ylabel={$f(t)$},
 title={同尺度对照：$t^{\log_2 3}$},title style={font=\small,yshift=3pt},
 tick label style={font=\small},label style={font=\small},
 ylabel style={at={(axis description cs:0,1)},anchor=south east,xshift=-18pt},
 clip=false,grid=major,grid style={black!8},enlargelimits=false,
 every axis plot/.append style={thick},samples=161]
\addplot[CExplanation,very thick,domain=.5:8.1] {x^(ln(3)/ln(2))};
\draw[dashed,black!35] (axis cs:2,0) |- (axis cs:0,3);
\draw[dashed,black!35] (axis cs:4,0) |- (axis cs:0,9);
\draw[dashed,black!35] (axis cs:8,0) |- (axis cs:0,27);
\addplot[only marks,mark=*,mark size=2pt,orange!80!black] coordinates {(2,3) (4,9) (8,27)};

\node[font=\small,anchor=north,yshift=-3pt] at (axis cs:2,-4) {$2$};
\node[font=\small,anchor=north,yshift=-3pt] at (axis cs:4,-4) {$4$};
\node[font=\small,anchor=north,yshift=-3pt] at (axis cs:8,-4) {$8$};

\end{axis}
\end{tikzpicture}
\end{minipage}
\par\small 都过 $(2,3)$，在 $4$ 处却读到 $6$ 和 $9$，在 $8$ 处读到 $9$ 和 $27$。一个点的值不能代替运算结构。
\end{center}

\end{examplebox}

\FReserve{4cm}
\begin{examplebox}
\textbf{\textcolor{CExample}{例 10（进阶）｜如何区别识别模型和确定解析式}}

下面都满足某种“正弦型”关系：
\[
f(x+y)f(x-y)=f(x)^2-f(y)^2.
\]
\textbf{函数一：} $f(x)=\sin x$，来自正弦平方差公式。

\textbf{函数二：} $f(x)=x$，因为
\[(x+y)(x-y)=x^2-y^2.\]
因此已知这条关系，不能就说 $f(x)=\sin x$。

\textbf{考试做法：}如果只求某两个函数值的乘积，直接代入关系式通常更快，不必尝试唯一确定 $f(x)$。
\end{examplebox}

